\documentclass[11pt,a4paper,oneside]{report}
\usepackage[spanish,activeacute,es-noshorthands]{babel}
\usepackage[utf8]{inputenc}
\usepackage[T1]{fontenc}
\usepackage{lmodern}
\usepackage{amsmath, amssymb}
\usepackage{graphicx}
\usepackage[dvipsnames]{xcolor}
\usepackage{geometry}
\geometry{margin=2.4cm, headheight=14pt}
\usepackage{microtype}
\usepackage{booktabs}
\usepackage{longtable}
\usepackage{enumitem}
\usepackage{fancyhdr}
\usepackage{tcolorbox}
\tcbuselibrary{skins, breakable}
\usepackage{caption}
\captionsetup{font=small, labelfont=bf}
\usepackage{titlesec}
\usepackage{hyperref}
\hypersetup{
  colorlinks=true,
  linkcolor=NavyBlue,
  citecolor=ForestGreen,
  urlcolor=BrickRed,
  pdftitle={Referencias open-source para CubeSat-EdgeAI-Payload},
  pdfauthor={INTISAT}
}

\definecolor{intisat}{RGB}{26, 54, 93}
\definecolor{intisataccent}{RGB}{197, 143, 0}
\definecolor{ok}{RGB}{47, 133, 90}
\definecolor{warn}{RGB}{197, 60, 60}
\definecolor{prio}{RGB}{197, 143, 0}

\titleformat{\chapter}[hang]
  {\normalfont\Huge\bfseries\color{intisat}}
  {\thechapter.}{0.6em}{}
\titlespacing*{\chapter}{0pt}{-20pt}{20pt}

\pagestyle{fancy}
\fancyhf{}
\fancyhead[L]{\nouppercase{\leftmark}}
\fancyhead[R]{\thepage}
\renewcommand{\headrulewidth}{0.4pt}

\newtcolorbox{recomendacion}[1][]{
  colback=yellow!7, colframe=intisataccent,
  fonttitle=\bfseries, title=Recomendacion,
  breakable, sharp corners=south, #1
}
\newtcolorbox{nota}[1][]{
  colback=blue!4, colframe=intisat,
  fonttitle=\bfseries, title=Como usarlo en tu proyecto,
  breakable, sharp corners=south, #1
}
\newtcolorbox{cuidado}[1][]{
  colback=red!4, colframe=warn,
  fonttitle=\bfseries, title=Cuidado,
  breakable, sharp corners=south, #1
}

\begin{document}

\begin{titlepage}
\centering
\vspace*{1cm}
{\color{intisat}\rule{\textwidth}{2pt}}
\vspace{0.6cm}

{\Huge\bfseries\color{intisat} Referencias open-source}\\[0.4em]
{\LARGE\bfseries para el desarrollo del CubeSat}

\vspace{0.7cm}
{\Large\itshape 20 repositorios y recursos curados\\
para construir flight software, ground station y simulacion}

\vspace{1cm}
{\color{intisataccent}\rule{0.6\textwidth}{1pt}}

\vspace{2cm}

\begin{tabular}{rl}
\textbf{Proyecto:} & CubeSat-EdgeAI-Payload \\
\textbf{Plataforma:} & INTISAT --- payload de microscopia FPM \\
\textbf{Audiencia:} & Equipo tecnico de 4 personas \\
\textbf{Plan asociado:} & 2 semanas de exploracion en paralelo \\
\textbf{Repositorio:} & \texttt{github.com/JaminYC/CubeSat-EdgeAI-Payload} \\
\textbf{Fecha:} & \today \\
\textbf{Documento:} & v1.0
\end{tabular}

\vfill

\begin{minipage}{0.85\textwidth}
\textbf{Resumen.}\quad
Compilacion curada de 20 repositorios open-source y recursos academicos
relevantes para el desarrollo de un CubeSat con IA a bordo. Cubre flight
software (NASA cFS, F'), ground station (OpenMCT), simulacion (42, GMAT),
autonomia (PLEXIL), CubeSats academicos completos (OreSat, UPSat, AAUSAT)
e IA en orbita (ESA Phi-Lab). Para cada referencia incluye URL,
descripcion, prioridad para el proyecto, y como aplicarla concretamente
al payload INTISAT. Acompaña al plan de exploracion paralela del equipo
de 4 personas durante 2 semanas.
\end{minipage}

\vspace{1.2cm}
{\color{intisat}\rule{\textwidth}{2pt}}
\end{titlepage}

\tableofcontents
\clearpage

% =====================================================================
\chapter{Como leer este documento}

Las 20 referencias estan ordenadas por categoria y dentro de cada
categoria por \textbf{prioridad para tu proyecto}. Cada entrada incluye:

\begin{itemize}[leftmargin=2em]
  \item \textbf{Nombre y URL} del repositorio o recurso.
  \item \textbf{Tipo}: framework / app / libro / paper / dashboard / etc.
  \item \textbf{Que es}: descripcion corta.
  \item \textbf{Como aplicalo al INTISAT}: caja azul.
  \item \textbf{Prioridad}: $\bigstar\bigstar\bigstar$ (top), $\bigstar\bigstar$ (medio),
        $\bigstar$ (futuro / nice-to-have).
  \item \textbf{Tiempo estimado de exploracion}.
\end{itemize}

\begin{recomendacion}
Si el equipo es de 4 personas y solo tienen 2 semanas, atacar primero
las marcadas con $\bigstar\bigstar\bigstar$. Cada miembro toma una
referencia distinta y la presenta al equipo al final.
\end{recomendacion}

% =====================================================================
\chapter{NASA --- el clasico (4 referencias)}

\section{1. NASA cFS --- Core Flight System}

\begin{itemize}[leftmargin=1.5em]
  \item URL: \url{https://github.com/nasa/cFS}
  \item Tipo: \textbf{Framework de flight software}
  \item Prioridad: $\bigstar\bigstar$ (medio --- arquitectura util)
  \item Tiempo: 5--10 dias
\end{itemize}

\textbf{Que es}: framework de NASA para flight software. Goddard lo usa
desde 2007 en TODAS las misiones. Diseñado para RTOS pero corre tambien
en Linux.

\textbf{Componentes clave}:
\begin{itemize}[leftmargin=2em]
  \item cFE (Core Flight Executive): runtime con publish/subscribe.
  \item OSAL: capa de abstraccion sobre Linux/RTEMS/FreeRTOS/VxWorks.
  \item PSP: adaptacion a hardware.
  \item Apps: limit checker (LC), memory manager (MM), file manager (FM),
        data storage (DS), housekeeping (HK), scheduler (SCH).
\end{itemize}

\begin{nota}
El \textbf{patron arquitectural} de cFS (publish/subscribe + tablas de
configuracion + FDIR) es directamente aplicable a tu daemon. NO copiar
codigo (es complejo), sino imitar la estructura: cada subsistema publica
mensajes, otros se suscriben, hay tablas de configuracion en tiempo de
ejecucion. Tu \texttt{cubesat/daemon.py} ya hace algo similar via
\texttt{/run/cubesat/}.
\end{nota}

\section{2. NASA F' (FPrime)}

\begin{itemize}[leftmargin=1.5em]
  \item URL: \url{https://github.com/nasa/fprime}
  \item Tipo: \textbf{Framework de flight software (moderno)}
  \item Prioridad: $\bigstar\bigstar\bigstar$ (alto)
  \item Tiempo: 3--7 dias
\end{itemize}

\textbf{Que es}: framework alternativo a cFS, mas reciente, mas simple,
mas accesible. Usado en misiones reales: \textbf{MarCO} (los CubeSats
interplanetarios que escoltaron a InSight) e \textbf{Ingenuity} (el
helicoptero de Mars).

\textbf{Por que vale}:
\begin{itemize}[leftmargin=2em]
  \item Mucho mejor documentado que cFS.
  \item Tutorial oficial paso a paso: \url{https://nasa.github.io/fprime/Tutorials/index.html}
  \item Generador de codigo a partir de XML.
  \item Hello World corre en 1 hora.
\end{itemize}

\begin{nota}
La opcion mas viable si quieren \textbf{adoptar un framework formal} en
el INTISAT. Puede coexistir con tu daemon Python: F' maneja la coordinacion
inter-subsistema y tu daemon corre el pipeline IA como un componente F'.
\end{nota}

\section{3. NASA OpenMCT}

\begin{itemize}[leftmargin=1.5em]
  \item URL: \url{https://github.com/nasa/openmct}
  \item Tipo: \textbf{Dashboard de mission control}
  \item Prioridad: $\bigstar\bigstar\bigstar$ (alto)
  \item Tiempo: 2 dias
\end{itemize}

\textbf{Que es}: framework web (JavaScript) para construir interfaces de
control de misiones. Lo usa JPL en operaciones reales. Permite integrar
telemetria de cualquier fuente y armar sinopticos, graficos, alarmas.

\textbf{Por que vale}:
\begin{itemize}[leftmargin=2em]
  \item Te ahorra meses de desarrollo de tu ground station.
  \item Tutorial en 30 minutos te deja un dashboard andando.
  \item Plugin architecture: conectas tu fuente de telemetria.
  \item Imprimible, exportable, theme-able.
\end{itemize}

\begin{nota}
\textbf{Una persona del equipo deberia probarlo esta semana}. Conectar
un script Python que lea \texttt{/run/cubesat/telemetry.json} (lo que ya
publica tu daemon) y servirlo a OpenMCT por WebSocket. En 2 horas tienen
el primer prototipo del dashboard del INTISAT.
\end{nota}

\section{4. NASA Astrobee}

\begin{itemize}[leftmargin=1.5em]
  \item URL: \url{https://github.com/nasa/astrobee}
  \item Tipo: \textbf{Software completo de un robot espacial}
  \item Prioridad: $\bigstar\bigstar$ (medio)
  \item Tiempo: 3 dias para entender, semanas para profundizar
\end{itemize}

\textbf{Que es}: el flight software de los 3 robots free-flying que vuelan
dentro de la ISS. Open source completo: vision, navegacion, planificacion,
comunicacion, control.

\begin{nota}
Util si quieren ver como NASA estructura un sistema con \textbf{vision
computacional + autonomia + telemetria} todo junto. Tu payload tiene la
misma combinacion (vision + IA + telemetria) en menor escala.
\end{nota}

% =====================================================================
\chapter{NASA --- apps puntuales de cFS (3 referencias)}

\section{5. cFS sample\_app}

\begin{itemize}[leftmargin=1.5em]
  \item URL: \url{https://github.com/nasa/sample_app}
  \item Tipo: \textbf{App de ejemplo de cFS}
  \item Prioridad: $\bigstar\bigstar\bigstar$ (alto si quieren entender cFS)
  \item Tiempo: 4--8 horas
\end{itemize}

\textbf{Que es}: el "Hello World" de cFS. Una app minima que muestra como
se estructura cualquier subsistema en el ecosistema cFS. Ideal para
aprender los conceptos sin perderse en el monstruo completo.

\begin{nota}
Antes de leer cFS completo, leer sample\_app. Te explica:
publish/subscribe, telemetria, comandos, tablas de configuracion. En 4 horas
ya tenes los conceptos.
\end{nota}

\section{6. cFS Limit Checker (LC)}

\begin{itemize}[leftmargin=1.5em]
  \item URL: \url{https://github.com/nasa/lc}
  \item Tipo: \textbf{App de FDIR}
  \item Prioridad: $\bigstar\bigstar$
  \item Tiempo: 1 dia
\end{itemize}

\textbf{Que es}: monitorea telemetria contra limites configurables y
dispara acciones de FDIR (Fault Detection, Isolation, Recovery) cuando
algo sale de rango. Ej: si la temperatura de la CPU supera 80°C, dispara
un comando de SAFE\_MODE.

\begin{nota}
Tu daemon hoy NO tiene FDIR formal. Esta app es la \textbf{plantilla de
referencia} para implementarlo. Conceptos: tabla de limites, watchpoints,
actionpoints, response actions.
\end{nota}

\section{7. cFS Scheduler (SCH)}

\begin{itemize}[leftmargin=1.5em]
  \item URL: \url{https://github.com/nasa/sch}
  \item Tipo: \textbf{Scheduler de mision}
  \item Prioridad: $\bigstar$ (futuro)
  \item Tiempo: 1 dia
\end{itemize}

\textbf{Que es}: ejecuta tareas periodicas (housekeeping, telemetria,
recolecciones cientificas) segun un schedule configurable. Ej: cada 30 s
publica HK, cada 1 hora hace un scan de payload, etc.

\begin{nota}
Tu sistema es REACTIVO (espera comandos del OBC), no SCHEDULED. Pero si
en algun momento queres un modo "autonomo cientifico" donde la payload
decida sola cuando capturar, el patron de SCH es el de referencia.
\end{nota}

% =====================================================================
\chapter{Simulacion y planeamiento (3 referencias)}

\section{8. NASA 42 simulator}

\begin{itemize}[leftmargin=1.5em]
  \item URL: \url{https://github.com/nasa/42}
  \item Tipo: \textbf{Simulador 6-DOF de spacecraft}
  \item Prioridad: $\bigstar\bigstar$ (alto si tienen ADCS team)
  \item Tiempo: 2--3 dias para arrancar
\end{itemize}

\textbf{Que es}: simulador completo de dinamica de naves espaciales. Calcula
orbita, actitud, sensores, actuadores, perturbaciones. Permite probar
algoritmos ADCS sin satelite real.

\begin{nota}
El equipo ADCS del INTISAT (no vos, el de actitud) puede probar sus
algoritmos aqui antes del integration test. Para tu payload no aplica
directamente.
\end{nota}

\section{9. NASA GMAT}

\begin{itemize}[leftmargin=1.5em]
  \item URL: \url{https://github.com/nasa/GMAT}
  \item Tipo: \textbf{General Mission Analysis Tool}
  \item Prioridad: $\bigstar\bigstar$ (alto para planeamiento)
  \item Tiempo: 1 dia para basicos
\end{itemize}

\textbf{Que es}: la herramienta de NASA para diseño y analisis de mision.
Calcula orbitas, eclipses, ground station passes, ventanas de comunicacion,
delta-v, tiempo de vida de bateria.

\begin{nota}
\textbf{Una persona del equipo de operaciones} deberia tener GMAT
configurado con la orbita real del INTISAT. Asi se predicen los pases
para downlink, los eclipses (en SAFE\_MODE) y se planifica el schedule
de scans.
\end{nota}

\section{10. NASA Trick}

\begin{itemize}[leftmargin=1.5em]
  \item URL: \url{https://github.com/nasa/trick}
  \item Tipo: \textbf{Framework de simulacion}
  \item Prioridad: $\bigstar$ (futuro)
  \item Tiempo: 1 semana
\end{itemize}

\textbf{Que es}: framework general de simulacion. Usado para crear
simuladores de sistemas (no solo spacecraft). NASA lo usa para entrenar
operadores y para testing pre-vuelo.

% =====================================================================
\chapter{Autonomia y vision (2 referencias)}

\section{11. NASA PLEXIL}

\begin{itemize}[leftmargin=1.5em]
  \item URL: \url{https://github.com/nasa/PLEXIL}
  \item Tipo: \textbf{Lenguaje y motor de autonomia}
  \item Prioridad: $\bigstar\bigstar$ (interesante para tu IA)
  \item Tiempo: 1 semana
\end{itemize}

\textbf{Que es}: Plan Execution Interchange Language. Define planes
ejecutables con condiciones, retries, recovery. Usado en rovers de Mars y
en autonomous fault recovery.

\begin{nota}
Si en algun momento quieren que el payload INTISAT \textbf{decida solo}
cuando hacer un scan, que mascara usar, que hacer si falla la primera, etc.,
PLEXIL es la herramienta. No la van a usar en la primera mision pero vale
saber que existe.
\end{nota}

\section{12. JPL Open Source Rover}

\begin{itemize}[leftmargin=1.5em]
  \item URL: \url{https://github.com/nasa-jpl/open-source-rover}
  \item Tipo: \textbf{Rover educativo}
  \item Prioridad: $\bigstar$ (recreativo / didactico)
  \item Tiempo: 1 dia para entender
\end{itemize}

\textbf{Que es}: rover construible en casa con piezas COTS, basado en
Curiosity y Perseverance. Documenta arquitectura mecanica + software de
movilidad + vision.

\begin{nota}
No directamente aplicable, pero \textbf{excelente para inspiracion}. El
software de control del rover (en GitHub) usa patrones similares a tu
daemon: comandos, telemetria, vision.
\end{nota}

% =====================================================================
\chapter{ESA y la IA en orbita (3 referencias)}

\section{13. ESA Phi-Lab}

\begin{itemize}[leftmargin=1.5em]
  \item URL: \url{https://github.com/ESA-PhiLab}
  \item Tipo: \textbf{Repositorio del laboratorio de IA de ESA}
  \item Prioridad: $\bigstar\bigstar\bigstar$ (top --- directamente alineado)
  \item Tiempo: 3--5 dias
\end{itemize}

\textbf{Que es}: la division de ESA dedicada a IA en orbita. Repos con codigo
de PhiSat-1, PhiSat-2, datasets, herramientas de inferencia para Movidius
Myriad 2, frameworks de evaluacion.

\textbf{Repos destacados}:
\begin{itemize}[leftmargin=2em]
  \item \texttt{ESA-PhiLab/PhiSat-1\_Coastline} --- modelo de deteccion de costas
  \item \texttt{ESA-PhiLab/Eyes-on-the-Earth} --- IA para EO en general
  \item \texttt{ESA-PhiLab/iris} --- pipeline de inferencia on-board
\end{itemize}

\begin{nota}
\textbf{Esto es lo mas cercano a lo que estan haciendo}. Leer los papers
de PhiSat-1 + el repo de \texttt{iris} les muestra:
\begin{itemize}[leftmargin=1em]
  \item Como deployar un modelo entrenado a hardware embebido en orbita.
  \item Que arquitecturas de red funcionan en Movidius / Pi.
  \item Como filtrar imagenes utiles vs no utiles en tiempo real.
\end{itemize}
\textbf{Lectura prioritaria del equipo}.
\end{nota}

\section{14. ESA Operations}

\begin{itemize}[leftmargin=1.5em]
  \item URL: \url{https://github.com/esa}
  \item Tipo: \textbf{Repositorios oficiales de ESA}
  \item Prioridad: $\bigstar\bigstar$
  \item Tiempo: navegar 1 dia
\end{itemize}

\textbf{Que es}: repositorio principal de ESA con varios subgrupos: ECSS
parsers, OPS-SAT experiments, MOOC materials, ESA herramientas.

\section{15. ESA OPS-SAT experiments}

\begin{itemize}[leftmargin=1.5em]
  \item URL: \url{https://nmf.esa.int/} y \url{https://github.com/esa/ops-sat-tutorials}
  \item Tipo: \textbf{Plataforma + tutoriales de IA en orbita}
  \item Prioridad: $\bigstar\bigstar\bigstar$
  \item Tiempo: 2--3 dias
\end{itemize}

\textbf{Que es}: OPS-SAT es el satelite de ESA para experimentar con IA y
software experimental in-orbit. Sus tutoriales muestran como subir un
experimento al satelite, correrlo, recibir resultados.

\begin{nota}
\textbf{Inspiracion directa}: si en futuras misiones del INTISAT quieren
correr experimentos cambiables, este es el modelo. Y si tu payload se
postula a OPS-SAT VOLT como guest experiment, conocer esta plataforma es
critico.
\end{nota}

% =====================================================================
\chapter{CubeSats academicos open-source (3 referencias)}

\section{16. OreSat (Portland State University)}

\begin{itemize}[leftmargin=1.5em]
  \item URL: \url{https://github.com/oresat}
  \item Tipo: \textbf{CubeSat 100\% open source}
  \item Prioridad: $\bigstar\bigstar\bigstar$ (top)
  \item Tiempo: 1 semana para profundizar
\end{itemize}

\textbf{Que es}: el CubeSat mas open de la historia. Hardware, firmware,
software, ground station, todo open. Diseñado por Portland State Aerospace
Society (PSAS). Han volado misiones reales.

\textbf{Componentes destacados}:
\begin{itemize}[leftmargin=2em]
  \item \texttt{oresat-c3} --- Command, Control, Communications software.
  \item \texttt{oresat-cans} --- comunicacion CAN entre subsistemas.
  \item \texttt{oresat-mainboard} --- el carrier board del OBC.
  \item \texttt{oresat-firmware} --- firmware multi-subsistema.
\end{itemize}

\begin{nota}
\textbf{Es el ejemplo mas cercano a tu situacion}: CubeSat academico, RPi
y MCUs, comunicacion por CAN, ground station propia. Antes de inventar
nada en el INTISAT, mirar como lo hicieron en OreSat.
\end{nota}

\section{17. UPSat (Libre Space Foundation)}

\begin{itemize}[leftmargin=1.5em]
  \item URL: \url{https://gitlab.com/librespacefoundation/upsat}
  \item Tipo: \textbf{CubeSat open-source europeo}
  \item Prioridad: $\bigstar\bigstar$
  \item Tiempo: 4 dias
\end{itemize}

\textbf{Que es}: primer CubeSat completamente open-source en volar. Lanzado
en 2017 desde la ISS. Toda la documentacion es publica: requirements,
diseño, codigo, tests.

\begin{nota}
Util para ver \textbf{como se documenta un CubeSat para revisores}. Si en
algun momento INTISAT pasa por una review formal (Critical Design Review,
Flight Readiness Review), tener UPSat de referencia ayuda mucho.
\end{nota}

\section{18. AAUSAT-3 / -4 (Aalborg University)}

\begin{itemize}[leftmargin=1.5em]
  \item URL: \url{http://www.space.aau.dk/}
  \item Tipo: \textbf{Documentacion academica}
  \item Prioridad: $\bigstar\bigstar$
  \item Tiempo: 2 dias
\end{itemize}

\textbf{Que es}: serie de CubeSats academicos de la Universidad de Aalborg
(Dinamarca). Documentacion publica con modos operativos detallados, FDIR,
state machines.

\begin{nota}
La \textbf{descripcion de modos de operacion} de AAUSAT es probablemente
la mejor del mundo academico. Si vas a diseñar el diagrama de estados
de mision del INTISAT, este es el template.
\end{nota}

% =====================================================================
\chapter{Estandares y handbooks (2 referencias)}

\section{19. CalPoly CubeSat Design Specification}

\begin{itemize}[leftmargin=1.5em]
  \item URL: \url{https://www.cubesat.org/cds-announcement}
  \item Tipo: \textbf{Estandar mecanico/electrico}
  \item Prioridad: $\bigstar\bigstar\bigstar$
  \item Tiempo: 4 horas
\end{itemize}

\textbf{Que es}: el estandar de facto de CubeSats. Define dimensiones,
masa, conector PC-104, requisitos de testing pre-vuelo. \textbf{Lectura
obligatoria} para cualquier CubeSat.

\section{20. NASA SE Handbook (Systems Engineering)}

\begin{itemize}[leftmargin=1.5em]
  \item URL: \url{https://www.nasa.gov/seh/}
  \item Tipo: \textbf{Manual de ingenieria de sistemas}
  \item Prioridad: $\bigstar\bigstar$
  \item Tiempo: 1 semana lectura selectiva
\end{itemize}

\textbf{Que es}: manual oficial de NASA. Define el proceso completo:
requirements, V\&V, reviews, modos operativos, FDIR, configuration
management.

\begin{nota}
No leer entero. Capitulos clave para CubeSat:
\begin{itemize}[leftmargin=1em]
  \item Cap. 4 --- Project Phases.
  \item Cap. 6 --- Operational Concept Document.
  \item Cap. 9 --- Verification \& Validation.
\end{itemize}
\end{nota}

% =====================================================================
\chapter{Plan de exploracion para el equipo de 4 personas}

Te propongo dividir el trabajo asi:

\begin{table}[h!]
\centering
\caption{Asignacion de exploracion --- 2 semanas.}
\renewcommand{\arraystretch}{1.4}
\begin{tabular}{p{1cm}p{4.5cm}p{4.5cm}p{4cm}}
\toprule
\textbf{Persona} & \textbf{Repos asignados} & \textbf{Lo que tiene que entender} & \textbf{Output esperado} \\
\midrule
A & OpenMCT (\#3) + ESA Phi-Lab (\#13) & Dashboard mission control + IA en orbita & Prototipo de UI conectado al daemon \\
B & cFS sample\_app (\#5) + LC (\#6) + F' (\#2) & Arquitectura pub/sub + FDIR & Doc 5 paginas con conceptos aplicables \\
C & OreSat (\#16) + AAUSAT (\#18) & CubeSat open completo + state machines & Inventario de subsistemas no documentados en INTISAT \\
D & 42 simulator (\#8) + GMAT (\#9) + UPSat (\#17) & Simulacion + planeamiento + docs formales & Borrador de orbita INTISAT con eclipses y pases \\
\bottomrule
\end{tabular}
\end{table}

\section{Cronograma}

\begin{itemize}[leftmargin=2em]
  \item \textbf{Dias 1--2}: cada persona instala las herramientas y corre los
        tutoriales basicos.
  \item \textbf{Dias 3--7}: profundizar en el repo asignado, tomar notas.
  \item \textbf{Dia 8}: reunion intermedia, brief de 10 min cada uno.
  \item \textbf{Dias 9--13}: profundizar en lo que se identifico como
        prioritario.
  \item \textbf{Dia 14}: reunion final --- decision sobre que 3 patrones se
        adoptan en INTISAT.
\end{itemize}

\section{Output del equipo al final de las 2 semanas}

\begin{enumerate}[leftmargin=2em]
  \item Decision sobre arquitectura: ¿adoptamos F'? ¿Mantenemos daemon
        Python con conceptos cFS? ¿Hibrido?
  \item Prototipo de OpenMCT funcionando con telemetria del payload.
  \item Borrador de modos de operacion de mision a nivel sistema.
  \item Plan de orbita INTISAT con eclipses y pases.
  \item Lista de patrones identificados en OreSat / AAUSAT que aplican.
\end{enumerate}

% =====================================================================
\chapter{Lectura priorizada (top 5 si solo tenes tiempo limitado)}

Si el equipo no puede dedicar 2 semanas, en orden:

\begin{enumerate}[leftmargin=2em]
  \item \textbf{OpenMCT tutorial} (\#3) --- 2 horas, salen ganando un dashboard.
  \item \textbf{cFS sample\_app README} (\#5) --- 4 horas, entiendes cFS sin perderte.
  \item \textbf{F' Tutorials} (\#2) --- 1 dia, alternativa moderna a cFS.
  \item \textbf{ESA Phi-Lab iris} (\#13) --- caso real de IA en orbita.
  \item \textbf{OreSat C3} (\#16) --- ver un CubeSat open completo.
\end{enumerate}

% =====================================================================
\chapter{Recursos extra que no son repositorios pero suman}

\section{NASA Software Catalog}

\url{https://software.nasa.gov}

Contiene codigo NASA NO publico en GitHub. Se descarga gratis con
registro. Hay herramientas valiosas (analisis termico, control,
simulacion) que no estan en GitHub.

\section{Small Satellite Conference USU --- Open Access Papers}

\url{https://digitalcommons.usu.edu/smallsat/}

\textbf{Todas} las papers de las conferencias SmallSat de los ultimos 30
años, gratis. El recurso \#1 academico de CubeSats. Buscar por keyword:
\textit{state machine}, \textit{FDIR}, \textit{operational modes},
\textit{onboard AI}.

\section{ESA SAVOIR architecture}

\url{https://savoir.estec.esa.int/}

Reference architecture de ESA para spacecraft on-board software.
Documentacion publica del esquema "como se hace un OBC en Europa".

\section{NASA cFS resource center}

\url{https://cfs.gsfc.nasa.gov/}

Documentacion central de cFS, tutorials, training material en video.

% =====================================================================
\chapter{Cierre}

Estas 20 referencias mas los 4 recursos extras te dan el universo de
lo que existe abierto en el ecosistema CubeSat. \textbf{No es para leerlas
todas}, es para que el equipo elija las 5 mas relevantes y las explore
con rigor.

\textbf{La leccion central}: NASA, ESA, y los proyectos academicos
abiertos (OreSat, UPSat, AAUSAT) llevan 30 años resolviendo los problemas
de CubeSat. Antes de inventar nada en el INTISAT, vale la pena ver como
ya se resolvio antes y adaptar.

\textbf{Lo que NO podes copiar} es la operacion: tu payload de microscopia
biologica con OLED programable y mascaras intercambiables es \textbf{nuevo}.
Eso es donde el equipo aporta valor diferencial. Pero la \textbf{arquitectura
de software, comunicacion, FDIR y operaciones} ya esta resuelta y
documentada. Hay que aprovecharla.

\end{document}
